\usepackage{amsthm}
\usepackage{amssymb}
\usepackage{nicematrix}

% \setmathfont[range={\doubleplus}, Scale=MatchLowercase]{Asana Math}

% Для отношения выводимости в грамматике. Смотри определение "Отношение выводимости".
\NewDocumentCommand{\derives}{O{*}}{\xRightarrow[]{#1}}

\newcommand{\Bbbzero}{\mathbb{0}}
\newcommand{\Bbbone}{\mathbb{1}}
\newcommand{\first}[1][1]{\text{\textsc{first}}_{#1}}
\newcommand{\follow}[1][1]{\text{\textsc{follow}}_{#1}}

% Количество ненулевых элементов (NNZ --- Number of Non-Zeros).
% Аргумент --- матрица или вектор, для которого считаем NNZ.
\newcommand*\NNZ[1]{\operatorname{NNZ}(#1)}

% Временная сложность умножения двух матриц над кольцом размера n × n (Matrix Multiplication).
\newcommand*\MM[1]{\operatorname{MM}(#1)}

% Временная сложность умножения двух булевых матриц размера n × n (Boolean Matrix Multiplication).
\newcommand*\BMM[1]{\operatorname{BMM}(#1)}

% Временная сложность умножения двух матриц размера n × n над указанным (полу)кольцом (Semiring Matrix Multiplication).
\newcommand*\SMM[2]{\operatorname{SMM}_{#1}(#2)}

% Временная сложность объединения двух булевых матриц (Boolean Matrix Union).
\newcommand*\BMU[1]{\operatorname{BMU}(#1)}

% Временная сложность пересечения двух булевых матриц (Boolean Matrix Intersection).
\newcommand*\BMI[1]{\operatorname{BMI}(#1)}

% Вершины графа, автомата
\newcommand*\circled[1]{\tikz[baseline=(char.base)]{
            \node[shape=circle,draw,inner sep=2pt] (char) {#1};}}

% Стиль для имён алгебраических структур
\newcommand\algebraic[1]{ \mathbb{#1}}

% Option<T>
\newcommand*\Opt[1]{\textit{Opt}_{#1}}

% Произведение Кронекера
\newcommand*\kron{\boxtimes}

% Произведение Кронекера, когда надо явно указать, откуда взялась поэлементная операция
\newcommand*\kronFrom[1]{ \kron^{\algebraic{#1}}}

% Произвольная бинарная операция
\newcommand*\binop{\circ}

% Бинарная операция с явным указанием алгебраической структуры из которой её взяли.
\newcommand*\binopFrom[1]{ \binop^{\algebraic{#1}}}

% Абстрактное сложение (в том числе, как поэлементная операция над матрицами и векторами)
\newcommand*\opAdd{\oplus}

% Абстрактное сложение (в том числе, как поэлементная операция над матрицами и векторами) с явным указанием алгебраической структуры из которой его взяли.
\newcommand*\opAddFrom[1]{ \opAdd^{\algebraic{#1}}}

% Абстрактное умножение
\newcommand*\opMult{\otimes}

% Абстрактное умножение с явным указанием алгебраической структуры из которой его взяли.
\newcommand*\opMultFrom[1]{ \opMult^{\algebraic{#1}}}

% Умножение двух матриц (или матрицы и вектора) с явным указанием алгебраической структуры из которой взяли поэлементные операции.
\newcommand*\mmultFrom[1]{ \opAdd \! \overset{\algebraic{#1}}{.} \! \opMult }

% Умножение двух матриц (или матрицы и вектора)
\newcommand*\mmult{ \mmultFrom{} }

\newenvironment{scaledalign}[4]
  {
    \begingroup
    #1
    \setlength\arraycolsep{#2}
    \renewcommand{\arraystretch}{#3}
    \begin{center}
    \begin{equation}
    \begin{aligned}
    #4
  }
  {
    \end{aligned}
    \end{equation}
    \end{center}
    \endgroup
  }
